\section{Tinjauan Pustaka}

\subsection{Model Generatif Musik Simbolik}

DeepBach \parencite{hadjeres2017} dirancang untuk menghasilkan \textit{chorale} empat suara dalam gaya Bach. Model memperkirakan distribusi nada dari konteks menggunakan jaringan yang mencakup LSTM dan melakukan generasi melalui \textit{pseudo-Gibbs sampling}. Pengguna dapat menetapkan batasan posisi seperti nada, ritme, atau kadens. Prosedur ini berbeda dari generasi sekuensial satu arah; perbedaan tersebut harus diperhitungkan dalam penjelasan model. Pengembangnya mengevaluasi DeepBach melalui uji diskriminasi daring ``\textit{Bach or Computer}'' dan menganalisis contoh keluaran. Analisis contoh tersebut mencatat kesalahan komposisi seperti oktaf sejajar, tanpa menghitung frekuensinya pada seluruh keluaran.

Coconet \parencite{huang2017} dikembangkan oleh peneliti MILA, Université de Montréal, dan Google Brain. Model ini menggunakan pendekatan distribusi simultan seluruh suara (\textit{orderless NADE}) yang memungkinkan pelengkapan suara secara fleksibel dari kondisi parsial apa pun. Naskah asli Coconet mengevaluasi model melalui \textit{log-likelihood} pada korpus \textit{chorale} Bach dan evaluasi oleh manusia.

NotaGen \parencite{wang2025} adalah model \textit{transformer} simbolik berskala besar yang dilatih pada lebih dari 1,6 juta karya musik dan \textit{fine-tuned} (disesuaikan lebih lanjut melalui pelatihan tambahan pada dataset spesifik) pada sekitar 9.000 komposisi klasik berkualitas tinggi. Para pengembangnya menyatakan bahwa model ini meningkatkan musikalitas keluaran dibandingkan model pembanding. Evaluasi dalam naskah aslinya menggunakan skor CLaMP 2 dan uji A/B subjektif. Naskah tersebut tidak melaporkan analisis \textit{voice leading} berbasis kaidah harmoni.

\subsection{Evaluasi Keluaran Model Musik Generatif}

Fang et al.\ \parencite{fang2020}, dalam ``\textit{Bach or Mock? A Grading Function for Chorales in the Style of J.S. Bach}'', mengembangkan fungsi penilaian menggunakan music21. Mereka membandingkan distribusi fitur musikal sebuah \textit{chorale} dengan distribusi fitur korpus Bach melalui jarak Wasserstein berbobot dan mengevaluasi keluaran model \textit{Transformer}. Salah satu fitur tersebut adalah distribusi kesalahan paralel, yaitu kemunculan unisono, kuint, dan oktaf sejajar. Skor yang lebih rendah menunjukkan kedekatan yang lebih tinggi menurut fungsi tersebut. Metrik ini berbeda dari indeks penalti pada penelitian ini; keduanya tidak boleh disebut sebagai rumus yang sama.

Huang et al.\ \parencite{huang2019} melaporkan penerapan Coconet dalam Bach Doodle dan menganalisis gerak paralel pada keluarannya. Dengan music21, mereka menghitung kuint dan oktaf sejajar pada 21,8 juta harmonisasi, yaitu rata-rata 0,365 dan 0,391 per birama. Frekuensi gerak paralel lebih tinggi ketika melodi masukan berada di luar distribusi data latih, dan lebih sedikit gerak paralel berkorelasi dengan umpan balik positif pengguna. Mereka juga menemukan 132 kuint sejajar dan 51 oktaf sejajar pada 382 \textit{chorale} Bach yang digunakan sebagai data latih dan validasi; banyak di antaranya dapat dimaklumi karena muncul pada batas frasa atau bersama nada non-akor. Temuan ini relevan bagi penelitian ini dalam dua hal: penghitungan gerak paralel dengan music21 dapat diterapkan pada keluaran model dalam jumlah besar, dan definisi operasional perlu menangani konteks yang membuat suatu gerak paralel dapat diterima. Analisis tersebut hanya mencakup Coconet serta kuint dan oktaf sejajar.

Yang dan Lerch \parencite{yang2020} membahas evaluasi model generatif musik dan mengusulkan pengukuran objektif berbasis distribusi fitur musikal. Retkowski et al.\ \parencite{Retkowski2024} mengembangkan ukuran jarak distribusi untuk musik simbolik. Ukuran distribusional menilai kedekatan kumpulan keluaran dengan data acuan; ukuran tersebut tidak menunjukkan lokasi pelanggaran kaidah pada partitur tertentu.

Le et al.\ \parencite{le2025} meninjau metode \textit{natural language processing} untuk generasi dan \textit{information retrieval} musik simbolik. Ji et al.\ \parencite{Ji2023} meninjau representasi, algoritme, dan metode evaluasi dalam generasi musik simbolik berbasis pembelajaran mendalam. Zhou et al.\ \parencite{Zhou2024} mengevaluasi kemampuan penalaran musik pada model bahasa besar (\textit{Large Language Models}). Studi Zhou et al.\ menguji model bahasa umum, bukan DeepBach, Coconet, atau NotaGen.

music21 \parencite{cuthbert2010} adalah pustaka Python sumber terbuka untuk musikologi berbasis komputer yang dikembangkan di MIT. Pustaka ini digunakan secara luas untuk analisis, manipulasi, dan evaluasi representasi musik simbolik. Fang et al.\ \parencite{fang2020} menggunakan music21 untuk membangun fungsi penilaian \textit{chorale}. Huang et al.\ \parencite{huang2019} menggunakan music21 untuk menghitung kuint dan oktaf sejajar pada keluaran Coconet. Penelitian ini menggunakan music21 untuk mengolah sampel MIDI. MusPy \parencite{dong2020} merupakan contoh pustaka Python lain untuk generasi dan evaluasi musik simbolik.

\subsection{Posisi Penelitian}

Penelitian ini berfokus pada perbandingan keluaran model menggunakan kaidah gerak suara yang dirumuskan dari Strube. Perbandingan mencakup pengukuran tiap kaidah serta perbedaan representasi dan kendali generasi pada DeepBach, Coconet, dan NotaGen. Dibandingkan dengan analisis Bach Doodle, penelitian ini menambahkan DeepBach dan NotaGen, kaidah resolusi nada penuntun, dan perbandingan antar-\textit{conditioning level}. Berbeda dari ukuran distribusional, pengukuran ini mencatat setiap penandaan beserta lokasi musikalnya sehingga dapat diperiksa kembali pada partitur.

\section{Landasan Teori}

\subsection{Musik Simbolik dan Model Generatif Simbolik}

Musik simbolik merepresentasikan nada, durasi, dan hubungan antar-suara dalam bentuk yang dapat dibaca secara struktural, seperti MIDI atau notasi ABC. Representasi ini memungkinkan pemeriksaan gerak suara setelah penguraian dan pemetaan suara diverifikasi. Penelitian dibatasi pada keluaran simbolik karena instrumen yang dikembangkan menggunakan peristiwa nada, bukan sinyal audio, sebagai masukan.

DeepBach mengisi koordinat nada berdasarkan konteks musikal dan \textit{pseudo-Gibbs sampling} \parencite{hadjeres2017}. Coconet menggunakan representasi \textit{piano roll}, memprediksi bagian partitur yang dihapus, dan melakukan pengisian ulang melalui \textit{blocked Gibbs sampling} \parencite{huang2017}. NotaGen menghasilkan notasi ABC dengan model \textit{Transformer} dan metadata periode, komponis, serta instrumentasi \parencite{wang2025}.

Ketiga antarmuka tersebut tidak menyediakan kendali yang identik. Penetapan sopran atau bas pada masukan DeepBach dan Coconet perlu diverifikasi melalui pelestarian nada yang ditetapkan. Metadata \texttt{Choir SATB} pada NotaGen meminta kategori instrumentasi; metadata itu belum membuktikan pelestarian melodi tertentu. Perbandingan harus menjelaskan kendali yang benar-benar tersedia dan dijalankan.

Data pelatihan dan spesialisasi tugas juga berbeda. Hasil perbandingan akan menggambarkan keluaran model pada konfigurasi yang diuji, dengan mempertimbangkan faktor arsitektur, data latih, representasi, dan prosedur generasi. Pemisahan pengaruh masing-masing faktor memerlukan desain yang lebih terkendali daripada rancangan tiga model yang ada.

\subsection{Teori Harmoni Fungsional dan Kaidah Gerak Suara Gustav Strube}

Landasan teori fokus penelitian ini tersusun dalam tiga tingkat. Teori harmoni fungsional menjadi teori utama (\textit{grand theory}). Kaidah gerak suara Gustav Strube menjadi teori yang menurunkan teori utama tersebut ke dalam aturan penulisan empat suara. Definisi operasional pada Bab III menerapkan kaidah tersebut pada musik simbolik.

Harmoni fungsional menjelaskan akor berdasarkan perannya dalam tonalitas, terutama fungsi tonika, dominan, dan subdominan, serta hubungan antar-akor yang membentuk progresi dan kadens \parencite{schoenberg1954}. Dalam kerangka ini, gerak setiap suara dari satu akor ke akor berikutnya mengikuti kaidah \textit{voice leading}. Gustav Strube dalam \textit{The Theory and Use of Chords} \parencite{strube1928} menguraikan kaidah tersebut dalam buku teks harmoni. Tiga kaidah yang menjadi dasar variabel terikat penelitian ini adalah sebagai berikut.

Pertama, larangan kuint sejajar (\textit{parallel fifths}): dua suara dilarang bergerak secara paralel dalam arah yang sama membentuk interval kuint sempurna di dua momen berurutan. Kaidah ini berkaitan dengan prinsip independensi suara \parencite{harrison2020}.

Kedua, larangan oktaf sejajar (\textit{parallel octaves}): dua suara dilarang bergerak secara paralel membentuk interval oktaf. Pelanggaran ini meniadakan kontribusi independen salah satu suara terhadap tekstur harmoni.

Ketiga, kewajiban resolusi nada penuntun (\textit{leading tone resolution}): nada ketujuh dari tangga nada, yang memiliki tegangan harmonik tertinggi, harus bergerak ke atas menuju tonika (nada pertama) ketika frase harmonis mengindikasikan resolusi \parencite{schoenberg1954}. Kegagalan resolusi ini menghasilkan kesan harmoni yang menggantung atau tidak tuntas.

Definisi operasional ketiga kaidah mencakup konteks, peristiwa yang layak diperiksa, dan pengecualian. Resolusi nada penuntun memerlukan perhatian pada fungsi nada dalam konteks harmonis. Anotasi contoh musikal menjadi acuan untuk membandingkan penandaan algoritmis dengan interpretasi gerak suara.

\section{Asumsi dan Hipotesis}

\subsection{Asumsi}

Penelitian ini bertolak dari dua anggapan dasar. Pertama, ketiga kaidah Strube dapat diterapkan pada tekstur empat suara bergaya \textit{chorale} yang menjadi keluaran target model. Kedua, keluaran simbolik yang lolos kontrol kualitas penguraian memuat informasi nada, durasi, dan suara yang cukup untuk memeriksa gerak antar-suara. Kedua anggapan ini diperiksa melalui kutipan halaman kaidah dari edisi Strube yang dirujuk dan kontrol kualitas sampel pada Bab III.

\subsection{Hipotesis}

Hipotesis adalah jawaban sementara terhadap masalah penelitian yang kebenarannya harus diuji secara empiris \parencite{rangkuti2016}. Pertanyaan penelitian pertama bersifat deskriptif sehingga dijawab dengan statistik deskriptif tanpa hipotesis. Pertanyaan kedua dan ketiga dijawab melalui hipotesis berikut. Setiap hipotesis diuji untuk masing-masing kaidah Strube secara terpisah, dengan tingkat kepatuhan sebagai variabel terikat.

\begin{enumerate}
    \item Hipotesis untuk pertanyaan penelitian kedua, pada setiap kondisi yang mendukung perbandingan: \\
    $H_0$: tidak terdapat perbedaan tingkat kepatuhan terhadap kaidah Strube antara \textit{output} DeepBach, Coconet, dan NotaGen. \\
    $H_1$: terdapat perbedaan tingkat kepatuhan terhadap kaidah Strube pada sedikitnya satu model.
    \item Hipotesis untuk pertanyaan penelitian ketiga, pada setiap model: \\
    $H_0$: \textit{conditioning level} tidak berpengaruh terhadap tingkat kepatuhan \textit{output} terhadap kaidah Strube. \\
    $H_1$: \textit{conditioning level} berpengaruh terhadap tingkat kepatuhan \textit{output} terhadap kaidah Strube.
\end{enumerate}

Kedua hipotesis dirumuskan tanpa arah karena sumber yang ditinjau belum memberikan dasar untuk memperkirakan model atau \textit{conditioning level} mana yang menghasilkan kepatuhan lebih tinggi. Penolakan $H_0$ pada hipotesis pertama menunjukkan perbedaan keluaran model pada konfigurasi yang diuji. Hasil tersebut tidak menunjukkan bahwa perbedaan disebabkan oleh arsitektur model semata. Pada hipotesis kedua, \textit{conditioning level} diberikan oleh peneliti pada model yang sama, sehingga perbedaan antar-\textit{conditioning level} dapat dibaca sebagai pengaruh kendali tersebut pada model dan checkpoint yang diuji.
